%==============================================================================
% CIC0201 - Seguranca Computacional
% Trabalho de Implementacao 1 - Cifra de Vigenere
% Formato: artigo IEEE de conferencia (duas colunas), sem folha de capa.
%
% Compilacao:  pdflatex relatorio.tex   (duas vezes, para as referencias)
% Requer a classe IEEEtran.cls (pacote texlive-publishers ou Overleaf).
%==============================================================================
\documentclass[conference]{IEEEtran}

\usepackage[utf8]{inputenc}
\usepackage[T1]{fontenc}
\usepackage[brazil]{babel}
\usepackage{amsmath}
\usepackage{graphicx}
\usepackage{booktabs}
\usepackage{listings}
\usepackage{xcolor}
\usepackage{url}

% --- Estilo dos trechos de codigo -------------------------------------------
\definecolor{codecomment}{rgb}{0.35,0.45,0.35}
\definecolor{codekeyword}{rgb}{0.10,0.25,0.55}
\lstset{
  basicstyle=\ttfamily\scriptsize,
  keywordstyle=\color{codekeyword}\bfseries,
  commentstyle=\color{codecomment}\itshape,
  numbers=none,
  breaklines=true,
  columns=fullflexible,
  frame=single,
  framesep=3pt,
  language=Python,
  % Permite acentos dentro dos listings sob pdflatex.
  literate=%
    {á}{{\'a}}1 {à}{{\`a}}1 {â}{{\^a}}1 {ã}{{\~a}}1
    {é}{{\'e}}1 {ê}{{\^e}}1 {í}{{\'i}}1 {ó}{{\'o}}1
    {ô}{{\^o}}1 {õ}{{\~o}}1 {ú}{{\'u}}1 {ü}{{\"u}}1
    {ç}{{\c c}}1 {Á}{{\'A}}1 {Ã}{{\~A}}1 {Ç}{{\c C}}1
    {É}{{\'E}}1 {Ó}{{\'O}}1
}

% --- Conteudo a preencher: use marcadores e remova-os antes da entrega -------
% (nenhum marcador pendente permanece neste arquivo)

\begin{document}

\title{Trabalho de Implementação 1:\\
Cifra de Vigenère --- Implementação e Criptoanálise}

\author{
\begin{tabular}{cc}
\begin{tabular}{c}
Rafael de Lima Pereira\\
Matrícula: 242043277
\end{tabular}
&
\begin{tabular}{c}
Pedro Freitas França\\
Matrícula: 232006028
\end{tabular}
\\[1em]
\begin{tabular}{c}
Rian Kallebe da Silva Lisboa\\
Matrícula: 242012000
\end{tabular}
&
\begin{tabular}{c}
Arthur Martins Pereira de Souza\\
Matrícula: 241004499
\end{tabular}
\end{tabular}\\[1em]
Departamento de Ciência da Computação\\
Universidade de Brasília
}

\maketitle

\begin{abstract}
Este relatório descreve a implementação de um cifrador/decifrador da cifra de
Vigenère e o desenvolvimento de um ataque de recuperação de chave baseado em
propriedades estatísticas da linguagem. São apresentadas as decisões de
projeto quanto ao tratamento do alfabeto, a metodologia de criptoanálise
adotada, os resultados obtidos sobre os criptogramas fornecidos e uma análise
das limitações da abordagem.
\end{abstract}

\begin{IEEEkeywords}
cifra de Vigenère, criptoanálise, índice de coincidência, análise de frequência
\end{IEEEkeywords}

%==============================================================================
\section{Introdução}
%==============================================================================
A cifra de Vigenère, descrita no século XVI, é uma cifra de substituição
polialfabética que durante cerca de três séculos foi considerada
indecifrável, o que lhe rendeu a alcunha de \textit{le chiffre
indéchiffrable}~\cite{stallings}. Sua segurança repousa na ideia de que uma mesma letra do
texto claro pode ser mapeada em letras diferentes do criptograma, dependendo
da posição em que ocorre, o que achata a distribuição de frequências que
denuncia as cifras monoalfabéticas.

Esse achatamento, contudo, é apenas parcial: a chave se repete ciclicamente e,
uma vez conhecido o seu comprimento, o criptograma se decompõe em vários
textos cifrados por cifra de César, cada um deles vulnerável à análise de
frequência clássica.

Este trabalho está dividido em duas partes complementares. A Parte~I consiste
na implementação de um cifrador e decifrador, discutida na
Seção~\ref{sec:implementacao}. A Parte~II consiste na criptoanálise de dois
criptogramas de chave desconhecida --- um em português e outro em inglês ---
tratada nas Seções~\ref{sec:ataque} e~\ref{sec:resultados}. O código-fonte
completo está disponível em~\cite{repo}.

%==============================================================================
\section{A Cifra de Vigenère}
\label{sec:cifra}
%==============================================================================
Seja $\Sigma = \{A, B, \dots, Z\}$ o alfabeto de $n = 26$ letras, e seja cada
letra identificada com um inteiro em $\{0, 1, \dots, 25\}$. Dada uma mensagem
$P = p_0 p_1 \dots p_{m-1}$ e uma chave $K = k_0 k_1 \dots k_{\ell-1}$ de
comprimento $\ell$, o criptograma $C = c_0 c_1 \dots c_{m-1}$ é obtido por

\begin{equation}
c_i = (p_i + k_{i \bmod \ell}) \bmod n,
\label{eq:cifracao}
\end{equation}

e a decifração, pela operação inversa,

\begin{equation}
p_i = (c_i - k_{i \bmod \ell}) \bmod n.
\label{eq:decifracao}
\end{equation}

A cifra pode ser vista, portanto, como a aplicação intercalada de $\ell$
cifras de César independentes: todas as posições $i$ congruentes entre si
módulo $\ell$ são deslocadas pelo mesmo valor $k_{i \bmod \ell}$. Essa
observação é simultaneamente a descrição do algoritmo e o ponto de partida do
ataque descrito na Seção~\ref{sec:ataque}.

Vale notar que para $\ell = 1$ a cifra degenera na cifra de César, e que para
$\ell \geq m$ com chave aleatória e de uso único obtém-se o cifrador de
Vernam, comprovadamente seguro. A fragilidade da construção está exatamente
na faixa intermediária, de uso prático: $\ell$ pequeno em relação a $m$ e
chave dotada de estrutura linguística.

%==============================================================================
\section{Decisões de Implementação (Parte I)}
\label{sec:implementacao}
%==============================================================================
A implementação foi feita em Python~3, sem qualquer biblioteca de
criptografia: as Equações~\eqref{eq:cifracao} e~\eqref{eq:decifracao} estão
codificadas diretamente no módulo \texttt{vigenere/cipher.py}. O projeto é
organizado em quatro partes: o núcleo da cifra (\texttt{cipher.py}), o ataque
(\texttt{attack.py}), a interface de linha de comando (\texttt{cli.py}) e a
suíte de testes automatizados (\texttt{tests/}).

\subsection{Núcleo da cifra}
O laço central percorre a mensagem caractere a caractere, mantendo um índice
próprio para a chave, e é reproduzido de forma resumida a seguir:

\begin{lstlisting}
for char in text:
    letter = fold(char)
    if letter is None:
        # Fora do alfabeto: mantido tal e qual
        # (sem avancar a chave) no modo "preserve",
        # descartado no modo "strict".
        if alphabet == "preserve":
            result.append(char)
        continue

    # C_i = (P_i + K_i) mod 26
    # P_i = (C_i - K_i) mod 26
    base = ord("A") if letter.isupper() else ord("a")
    shift = offsets[key_index % len(offsets)] * sign
    shifted = chr((ord(letter) - base + shift)
                  % ALPHABET_SIZE + base)
    result.append(shifted)
    key_index += 1
\end{lstlisting}

O parâmetro \texttt{sign} vale $+1$ na cifração e $-1$ na decifração, de modo
que uma única rotina implementa as duas equações. Como \texttt{key\_index} é
incrementado apenas quando um caractere é efetivamente cifrado, a chave nunca
perde a sincronia com as letras às quais se aplica --- propriedade
indispensável para que a decifração seja o inverso exato da cifração.

\subsection{Tratamento do alfabeto}
O alfabeto adotado é o clássico de 26 letras. A decisão de projeto mais
relevante diz respeito aos caracteres que não pertencem a ele. A
Tabela~\ref{tab:alfabeto} resume o tratamento aplicado a cada classe de
caractere.

\begin{table}[htbp]
\caption{Tratamento de cada classe de caractere}
\label{tab:alfabeto}
\centering
\footnotesize
\begin{tabular}{@{}p{0.30\columnwidth}p{0.60\columnwidth}@{}}
\toprule
\textbf{Entrada} & \textbf{Tratamento} \\
\midrule
\texttt{a-z} / \texttt{A-Z} &
Cifradas. Caixa preservada no modo \texttt{preserve}; forçada a maiúscula no
modo \texttt{strict}. \\
\addlinespace
Letras acentuadas
(\texttt{á à â ã ç é ê í ó ô õ ú ü ñ}) &
Reduzidas à letra base por decomposição Unicode NFD e então cifradas. O acento
não é restaurado na decifração. \\
\addlinespace
Espaços, dígitos, pontuação &
\texttt{preserve}: copiados sem alteração, sem consumir letra da chave.
\texttt{strict}: descartados. \\
\addlinespace
Caracteres da chave &
Reduzidos da mesma forma; o que não é letra é ignorado. Chave sem nenhuma
letra é erro. \\
\bottomrule
\end{tabular}
\end{table}

\subsubsection{Redução de acentos por NFD}
Textos em português contêm letras acentuadas que não pertencem a $\Sigma$.
Duas alternativas foram consideradas: preservá-las intactas no criptograma, ou
reduzi-las à letra base. Optou-se pela segunda, implementada pela
normalização Unicode NFD, que decompõe \texttt{ç} em \texttt{c} seguido de uma
cedilha combinante, bastando então tomar o primeiro ponto de código:

\begin{lstlisting}
def fold(char):
    base = unicodedata.normalize("NFD", char)[:1]
    if "a" <= base <= "z" or "A" <= base <= "Z":
        return base
    return None
\end{lstlisting}

A justificativa é criptográfica. Preservar os acentos faria com que
\texttt{ç}, \texttt{ã} e \texttt{é} aparecessem em claro no criptograma, o que
não apenas vaza informação sobre o texto original como praticamente identifica
o idioma e restringe fortemente as palavras possíveis. A redução é ainda mais
robusta que uma lista fixa de caracteres acentuados: como opera sobre a
decomposição Unicode, trata corretamente qualquer letra latina acentuada,
inclusive as que não ocorrem em português. O custo dessa escolha é a perda de
informação: \texttt{ação} é recuperado como \texttt{acao}. Trata-se de perda
ortográfica, e não semântica, e é o comportamento usual em criptoanálise
clássica.

\subsubsection{Modos de alfabeto}
Foram implementados dois modos, selecionáveis pela opção
\texttt{-{}-alphabet}:

\begin{itemize}
  \item \texttt{preserve} (padrão): espaços, dígitos e pontuação são copiados
        sem alteração e \emph{não} consomem letra da chave; a caixa das letras
        é mantida. O resultado permanece legível, mas a estrutura do texto
        claro (limites de palavras e pontuação) continua visível no
        criptograma.
  \item \texttt{strict}: tudo o que não é letra é descartado e a saída é uma
        sequência contínua de letras maiúsculas. É o formato canônico dos
        exercícios de criptoanálise, e é o modo empregado na Parte~II, pois
        não revela limites de palavras.
\end{itemize}

A existência do modo \texttt{strict} tem, além disso, papel metodológico: ela
permite gerar criptogramas no mesmo formato dos fornecidos pelo enunciado, o
que possibilita validar o ataque contra textos de chave conhecida antes de
aplicá-lo aos criptogramas do trabalho.

\subsection{Validação das entradas}
A única condição que impede a execução é uma chave sem nenhuma letra do
alfabeto, que provoca erro explícito. As demais situações suspeitas geram
avisos na saída de erro padrão, sem interromper o processamento:

\begin{itemize}
  \item chave da qual foram removidos caracteres ou cujos acentos foram
        reduzidos --- é exibida a chave efetivamente utilizada, evitando que o
        usuário decifre com chave diferente da que julga estar usando;
  \item chave de uma única letra, equivalente à cifra de César;
  \item texto sem nenhuma letra do alfabeto, caso em que nada é cifrado.
\end{itemize}

A opção \texttt{-v} exibe um resumo dos parâmetros efetivos (comando,
alfabeto, chave utilizada e seu tamanho, número de caracteres e de letras da
entrada), e as opções \texttt{-i}/\texttt{-o} permitem ler e escrever
arquivos, necessárias para tratar os criptogramas da Parte~II.

\subsection{Verificação}
A implementação é acompanhada de uma suíte de 26 testes automatizados
(\texttt{pytest}) que cobrem: o vetor clássico
\texttt{ATTACKATDAWN}/\texttt{LEMON} $\rightarrow$ \texttt{LXFOPVEFRNHR};
a inversibilidade da cifra; a preservação da caixa; a redução de acentos e a
consequente sincronia da chave; o comportamento dos dois modos de alfabeto; o
tratamento de chaves com caracteres inválidos; e o comportamento da interface
de linha de comando, incluindo os avisos e as mensagens de erro.

%==============================================================================
\section{Metodologia do Ataque (Parte II)}
\label{sec:ataque}
%==============================================================================
A quebra de uma cifra polialfabética como a de Vigenère baseia-se na exploração
das propriedades estatísticas das linguagens naturais. Diferentemente de
sequências puramente aleatórias, idiomas como o português e o inglês possuem
distribuições de frequência de letras desiguais e previsíveis. O ataque explora
ainda o fato, observado na Seção~\ref{sec:cifra}, de que a cifra é a
intercalação de $\ell$ cifras de César. A abordagem fundamenta-se em duas
métricas --- o índice de coincidência (IC) e o teste de aderência
qui-quadrado ($\chi^2$) --- e é conduzida em duas etapas: primeiro estima-se
$\ell$, depois recupera-se cada letra da chave independentemente. A
implementação foi feita de forma modular em Python, reutilizando o decifrador
da Parte~I.

\subsection{Índice de coincidência}
O IC, introduzido por Friedman~\cite{friedman}, é a probabilidade de que dois
elementos selecionados aleatoriamente de um texto sejam idênticos. Para um
texto de $N$ letras com $n_i$ ocorrências da $i$-ésima letra do alfabeto,

\begin{equation}
\mathrm{IC} = \frac{\sum_{i=A}^{Z} n_i (n_i - 1)}{N (N-1)}.
\label{eq:ic}
\end{equation}

Para um texto aleatório uniforme o valor esperado é $1/26 \approx 0{,}038$;
para textos em linguagem natural situa-se em torno de $0{,}066$ em inglês e
$0{,}074$ em português, por causa da distribuição desigual das letras.
Essa desigualdade é o fundamento estatístico de todo o ataque: a
Tabela~\ref{tab:freq} apresenta as frequências relativas adotadas na
implementação~\cite{wikipedia_freq}, as mesmas usadas pelo teste
$\chi^2$ na recuperação da chave.

\begin{table}[htbp]
\caption{Frequência relativa das letras (\%) em português e em inglês}
\label{tab:freq}
\centering
\footnotesize
\begin{tabular}{@{}cccccc@{}}
\toprule
\textbf{Letra} & \textbf{PT} & \textbf{EN} &
\textbf{Letra} & \textbf{PT} & \textbf{EN} \\
\midrule
A & $14{,}63$ & $8{,}17$ & N & $5{,}05$  & $6{,}75$ \\
B & $1{,}04$  & $1{,}49$ & O & $10{,}73$ & $7{,}51$ \\
C & $3{,}88$  & $2{,}78$ & P & $2{,}52$  & $1{,}93$ \\
D & $4{,}99$  & $4{,}25$ & Q & $1{,}20$  & $0{,}10$ \\
E & $12{,}57$ & $12{,}70$ & R & $6{,}53$  & $5{,}99$ \\
F & $1{,}02$  & $2{,}23$ & S & $7{,}81$  & $6{,}33$ \\
G & $1{,}30$  & $2{,}02$ & T & $4{,}34$  & $9{,}06$ \\
H & $1{,}28$  & $6{,}09$ & U & $4{,}63$  & $2{,}76$ \\
I & $6{,}18$  & $6{,}97$ & V & $1{,}67$  & $0{,}98$ \\
J & $0{,}40$  & $0{,}15$ & W & $0{,}01$  & $2{,}36$ \\
K & $0{,}02$  & $0{,}77$ & X & $0{,}21$  & $0{,}15$ \\
L & $2{,}78$  & $4{,}03$ & Y & $0{,}01$  & $1{,}98$ \\
M & $4{,}74$  & $2{,}41$ & Z & $0{,}47$  & $0{,}07$ \\
\bottomrule
\end{tabular}
\end{table}

Observa-se, por exemplo, que A e E dominam o português ($14{,}63\%$ e
$12{,}57\%$), enquanto no inglês E e T são as mais frequentes ($12{,}70\%$ e
$9{,}06\%$). Letras como K, W e Y são praticamente ausentes em português.
Essa assimetria é o que o índice de coincidência detecta globalmente e o que
o qui-quadrado explorará localmente em cada coluna.

\subsection{Estimativa do comprimento da chave}
A estimativa é automatizada iterando sobre comprimentos candidatos
$k = 1, 2, \dots, 20$. Para cada $k$, o criptograma é subdividido em $k$
colunas --- cada coluna contém os caracteres cifrados pela mesma posição da
chave --- e calcula-se a média do IC dessas colunas. Quando $k$ é o
comprimento correto (ou um múltiplo dele), cada coluna resulta de uma única
cifra de César, que é uma permutação do alfabeto e portanto \emph{preserva} o
IC do idioma original: o IC médio aproxima-se de $0{,}07$. Quando $k$ está
errado, as colunas misturam deslocamentos distintos e o IC médio permanece
próximo de $0{,}038$. Seleciona-se o $k$ que maximiza o IC médio, pois indica
que os caracteres de cada coluna foram deslocados pelo mesmo caractere da
chave.

\subsection{Análise de frequência e recuperação da chave}
Fixado $k$, o texto permanece particionado em seus $k$ subconjuntos. Para cada
coluna, aplica-se força bruta sobre os $26$ deslocamentos possíveis da cifra de
César: a coluna é decifrada temporariamente e a distribuição resultante é
comparada à distribuição esperada do idioma (português ou inglês) pela
estatística qui-quadrado

\begin{equation}
\chi^2 = \sum_{i=A}^{Z} \frac{(O_i - E_i)^2}{E_i},
\label{eq:chi2}
\end{equation}

em que $O_i$ é o número de ocorrências observadas da letra $i$ e $E_i$ é o
número esperado, calculado a partir do tamanho da coluna e da probabilidade
teórica da letra no idioma (Tabela~\ref{tab:freq}). O deslocamento que
minimiza $\chi^2$ --- a menor divergência em relação ao idioma alvo --- é o
candidato mais provável para aquela posição da chave. A concatenação desses
caracteres reconstrói a chave completa e permite decifrar o criptograma.

\subsection{Avaliação e refinamento}
A chave candidata é usada para decifrar o criptograma, e o resultado é
avaliado pelo IC do texto decifrado --- que deve aproximar-se do valor
esperado do idioma --- e por inspeção da legibilidade. Quando alguma coluna é
curta demais para que a estatística seja confiável, o segundo e o terceiro
melhores candidatos de $\chi^2$ daquela posição são examinados, refinando-se a
chave posição a posição.

%==============================================================================
\section{Resultados}
\label{sec:resultados}
%==============================================================================
\subsection{Parte I}
O cifrador reproduz o vetor de teste clássico da literatura e é inversível
sobre texto em português, como ilustrado a seguir:

\begin{lstlisting}[language={}]
$ vigenere encode "ATTACKATDAWN" --key LEMON
LXFOPVEFRNHR
$ vigenere decode "LXFOPVEFRNHR" --key LEMON
ATTACKATDAWN
$ vigenere encode "Ataque ao amanhecer!" \
      --key segredo --alphabet strict
SXGHYHOGESRRKSUIX
$ vigenere decode "SXGHYHOGESRRKSUIX" \
      --key segredo --alphabet strict
ATAQUEAOAMANHECER
\end{lstlisting}

Note-se, no último exemplo, o efeito documentado do modo \texttt{strict}: a
pontuação e os espaços não são recuperados, pois foram descartados na
cifração. A suíte de 26 testes automatizados é executada integralmente com
sucesso.

\subsection{Parte II}
Dois criptogramas de validação --- um em português ($762$ letras) e outro em
inglês ($697$ letras), cifrados no modo \texttt{strict} com chaves distintas
--- foram processados pelo comando \texttt{vigenere attack}. O ataque
reproduz, na prática, as duas etapas da Seção~\ref{sec:ataque}: primeiro
estima o comprimento da chave pelo IC médio das colunas; depois recupera cada
letra pelo mínimo de $\chi^2$ frente às frequências da
Tabela~\ref{tab:freq}.

\subsubsection{Estimativa do comprimento da chave pelo IC}
Para cada criptograma, o módulo percorreu $k \in \{1,\dots,20\}$, particionou
o texto em $k$ colunas e calculou o IC médio. A Tabela~\ref{tab:ic} destaca
os comprimentos mais informativos. No criptograma em português, o IC médio
salta de valores próximos de $0{,}044$ (próximos do aleatório) para
$0{,}078$ em $k=9$, valor compatível com linguagem natural; o segundo pico
aparece em $k=18$ ($0{,}077$), múltiplo do comprimento correto. No inglês, o
máximo ocorre em $k=6$ ($0{,}068$), com picos secundários em $k=12$ e
$k=18$. Em ambos os casos o algoritmo seleciona automaticamente o $k$ de
maior IC médio --- $9$ e $6$, respectivamente --- sem intervenção manual.

\begin{table}[htbp]
\caption{IC médio das colunas para comprimentos selecionados}
\label{tab:ic}
\centering
\footnotesize
\begin{tabular}{@{}lccccc@{}}
\toprule
\textbf{Criptograma} & $k{=}1$ & $k{=}6$ & $k{=}9$ & $k{=}12$ & $k{=}18$ \\
\midrule
Português & $0{,}044$ & $0{,}052$ & $\mathbf{0{,}078}$ & $0{,}052$ & $0{,}077$ \\
Inglês    & $0{,}044$ & $\mathbf{0{,}068}$ & $0{,}053$ & $0{,}068$ & $0{,}068$ \\
\bottomrule
\end{tabular}
\end{table}

Esse comportamento confirma a intuição teórica: só quando $k$ coincide com o
período da chave (ou com um múltiplo) cada coluna resulta de uma única cifra
de César e o IC médio se eleva ao patamar do idioma; caso contrário, as
colunas misturam deslocamentos distintos e o IC permanece baixo.

\subsubsection{Recuperação da chave pela análise de frequência}
Fixado $k$, cada coluna é tratada como um César independente. Para cada uma
das $26$ hipóteses de deslocamento, a coluna é temporariamente decifrada e a
distribuição resultante é comparada à da Tabela~\ref{tab:freq} via
$\chi^2$~\eqref{eq:chi2}. O deslocamento de menor $\chi^2$ é a letra candidata
daquela posição. A Tabela~\ref{tab:chi2} ilustra os três melhores candidatos
em posições representativas: a margem entre o primeiro e o segundo lugar é
ampla (por exemplo, $\chi^2 \approx 22$ contra $\approx 749$ na primeira
posição do português), o que torna a escolha inequívoca.

\begin{table}[htbp]
\caption{Candidatos de $\chi^2$ por posição (letra:~valor; menor é melhor)}
\label{tab:chi2}
\centering
\footnotesize
\begin{tabular}{@{}cll@{}}
\toprule
\textbf{Pos.} & \textbf{PT ($k{=}9$)} & \textbf{EN ($k{=}6$)} \\
\midrule
$0$ & \texttt{S}:$22{,}1$; \texttt{T}:$749$; \texttt{F}:$758$
    & \texttt{C}:$25{,}7$; \texttt{P}:$374$; \texttt{I}:$386$ \\
$1$ & \texttt{E}:$14{,}0$; \texttt{S}:$813$; \texttt{P}:$919$
    & \texttt{I}:$20{,}6$; \texttt{B}:$542$; \texttt{O}:$579$ \\
$2$ & \texttt{G}:$21{,}7$; \texttt{T}:$204$; \texttt{H}:$398$
    & \texttt{P}:$32{,}0$; \texttt{F}:$422$; \texttt{A}:$473$ \\
$3$ & \texttt{U}:$17{,}5$; \texttt{F}:$514$; \texttt{H}:$694$
    & \texttt{H}:$49{,}8$; \texttt{U}:$237$; \texttt{A}:$250$ \\
$4$ & \texttt{R}:$26{,}5$; \texttt{C}:$547$; \texttt{Q}:$725$
    & \texttt{E}:$25{,}3$; \texttt{R}:$275$; \texttt{K}:$343$ \\
$5$ & \texttt{A}:$24{,}4$; \texttt{L}:$402$; \texttt{N}:$802$
    & \texttt{R}:$20{,}8$; \texttt{E}:$244$; \texttt{X}:$342$ \\
\bottomrule
\end{tabular}
\end{table}

Concatenando o melhor candidato de cada posição obtém-se
\texttt{SEGURANCA} (português) e \texttt{CIPHER} (inglês). Em nenhuma
posição foi necessário recorrer ao segundo ou ao terceiro colocado: o mínimo
de $\chi^2$ coincidiu com a letra correta em todas as colunas. A
Tabela~\ref{tab:resultados} resume o desfecho. Após a decifração, o IC do
texto claro aproximou-se do valor esperado do idioma ($0{,}079$ frente a
$\approx 0{,}074$ em português; $0{,}067$ frente a $\approx 0{,}066$ em
inglês).

\begin{table}[htbp]
\caption{Resultados da recuperação de chave}
\label{tab:resultados}
\centering
\footnotesize
\begin{tabular}{@{}lccc@{}}
\toprule
\textbf{Criptograma} & \textbf{$k$ estimado} & \textbf{Chave recuperada} &
\textbf{IC final} \\
\midrule
Português & $9$ & \texttt{SEGURANCA} & $0{,}079$ \\
Inglês    & $6$ & \texttt{CIPHER}    & $0{,}067$ \\
\bottomrule
\end{tabular}
\end{table}

%==============================================================================
\section{Limitações da Abordagem}
\label{sec:limitacoes}
%==============================================================================
O ataque descrito é estatístico e, como tal, depende de premissas que nem
sempre se verificam.

\textbf{Volume de texto.} Tanto o IC quanto o qui-quadrado são estimadores
que convergem lentamente. Cada subconjunto tem cerca de $N/k$ letras;
quando essa razão cai abaixo de algumas dezenas de caracteres, a distribuição
observada torna-se ruidosa e a letra correta da chave pode não minimizar
$\chi^2$. O ataque degrada, portanto, para criptogramas curtos ou chaves
longas, e falha por completo no caso limite $k \geq N$.

\textbf{Dependência do idioma.} O método pressupõe conhecimento prévio da
distribuição de frequências do idioma do texto claro. Textos técnicos, com
vocabulário atípico, ou textos que misturam idiomas produzem distribuições
distintas das tabelas de referência. Neste trabalho o idioma é informado na
execução do ataque (\texttt{-{}-language pt|en}); num cenário real seria
necessário testar as duas hipóteses e escolher a de melhor ajuste.

\textbf{Efeito do pré-processamento.} A redução de acentos altera levemente as
frequências do português --- ocorrências de \texttt{ç} são somadas às de
\texttt{c}, e as vogais acentuadas às suas bases. Como as tabelas de
referência usuais já são construídas sobre o alfabeto não acentuado, o efeito
é pequeno, mas introduz um viés que não é nulo.

\textbf{Múltiplos do comprimento.} A estimativa por IC não distingue $k$ de
seus múltiplos, já que particionar por $2k$ também isola cifras de César e
pode maximizar (ou empatar) o IC médio. Nos casos observados, a inspeção dos
valores de IC ao longo de $k = 1,\dots,20$ e a legibilidade do texto decifrado
permitem discriminar o comprimento verdadeiro do múltiplo.

\textbf{Escopo da implementação.} O cifrador opera sobre um alfabeto de 26
letras: dígitos e pontuação não são cifrados, apenas preservados ou
descartados. Uma mensagem cujo conteúdo relevante seja numérico não é
protegida pela implementação.

%==============================================================================
\section{Conclusão}
%==============================================================================
A implementação da cifra de Vigenère é direta: as decisões críticas estão no
alfabeto e no tratamento dos caracteres fora dele, com consequências
criptográficas concretas. A criptoanálise confirma que a repetição periódica
da chave é o ponto de ruptura --- permite decompor o problema em subcifras
independentes, de modo que a segurança depende da razão entre o comprimento
da chave e o do texto, não do tamanho do espaço de chaves.

%==============================================================================
\begin{thebibliography}{9}

\bibitem{stallings}
W. Stallings, \textit{Cryptography and Network Security: Principles and
Practice}, 7th ed. Upper Saddle River, NJ, USA: Pearson, 2017.

\bibitem{wikipedia_freq}
Wikipédia, ``Frequência de letras.'' [Online]. Disponível:
\url{https://pt.wikipedia.org/wiki/Frequência_de_letras}

\bibitem{friedman}
W. F. Friedman, \textit{The Index of Coincidence and Its Applications in
Cryptography}. Geneva, IL, USA: Riverbank Laboratories, 1922.

\bibitem{repo}
R. Lima, ``Implementação da cifra de Vigenère.'' [Online]. Disponível:
\url{https://github.com/RafaelLime/vigenere}

\end{thebibliography}

\end{document}
